% Part: lambda-calculus
% Chapter: introduction
% Section: syntax

\documentclass[../../../include/open-logic-section]{subfiles}

\begin{document}

\olfileid[fr]{lam}{int}{syn}
\olsection{La syntaxe du lambda-calcul}

Partons d'une suite de variables $x$, $y$,
$z$,~\dots{} et de symboles de constante $a$,
$b$, $c$,~\dots. L'ensemble des termes se
définit par induction :
\begin{enumerate}
\item Chaque variable est un terme.
\item Chaque constante est un terme.
\item Si $M$ et $N$ sont des termes, $(MN)$
      est un terme.
\item Si $M$ est un terme et $x$ une variable,
      $(\lambd[x][M])$ est un terme.
\end{enumerate}
Les termes de la forme $(MN)$ sont des
\emph{applications}; ceux de la forme
$(\lambd[x][M])$ sont des \emph{abstractions}.

Le système dépourvu de constantes s'appelle
le \emph{lambda-calcul pur}. Nous travaillerons
surtout dans ce système; toutes les lettres
minuscules y désigneront donc des variables.
Les lettres majuscules ($M$, $N$, etc.)
désigneront des termes du lambda-calcul.

Nous adopterons quelques conventions de
notation :

\begin{conv}
\begin{enumerate}
\item Lorsque les parenthèses sont omises,
      l'application s'associe de gauche à droite.
      Par exemple, si $M$, $N$, $P$ et $Q$ sont
      des termes, $MNPQ$ abrège $(((MN)P)Q)$.
\item Lorsque les parenthèses sont omises,
      l'abstraction lambda reçoit la portée la
      plus large possible. Ainsi,
      $\lambd[x][MNP]$ se lit
      $(\lambd[x][((MN)P)])$.
\item Un lambda peut abstraire plusieurs
      variables. Par exemple,
      $\lambd[xyz][M]$ abrège
      $\lambd[x][\lambd[y][\lambd[z][M]]]$.
\end{enumerate}
\end{conv}

Par exemple,
\[
\lambd[xy][xxyx \lambd[z][xz]]
\]
abrège
\[
\lambd[x][\lambd[y][((((xx)y)x)(\lambd[z][(xz)]))]].
\]
Retenez ces conventions. Elles risquent de
vous rendre fou au début; vous vous y
habituerez, et elles finiront par vous
agacer moins qu'un enchevêtrement de
parenthèses.

Deux termes qui ne diffèrent que par un renommage sans capture
de leurs variables liées sont dits
$\alpha$-équivalents; c'est le cas de
$\lambd[x][x]$ et $\lambd[y][y]$. Il sera
commode de les considérer comme le « même »
terme : lorsque nous dirons que $M$ et $N$
sont identiques, nous l'entendrons aussi à
un renommage des variables liées près.
Une occurrence de variable est \emph{liée} si elle se trouve
dans le corps d'une abstraction qui porte son nom.
Son lieur est alors l'abstraction de ce nom la plus proche.
Toute occurrence de variable qui n'a pas de lieur est \emph{libre}.
Les termes $\lambd[x][x]$ et $\lambd[y][y]$ n'ont aucune
variable libre. Dans
\[
(\lambd[z][yz])x,
\]
l'occurrence de $z$ dans le corps est liée par $\lambd[z]$.
Les occurrences de $y$ et de $x$ sont libres.

\end{document}
